% Finds the package and the example diagrams, also when compiled directly in an editor:
\makeatletter\def\input@path{{../}{../../app/samples/}}\makeatother
\documentclass[10pt,a4paper]{article}
\usepackage[a4paper,margin=2cm]{geometry}
\usepackage[T1]{fontenc}
\usepackage{lmodern}
\usepackage{gsmdiagram}
\usepackage[most]{tcolorbox}
\tcbuselibrary{listings}
\usepackage{booktabs,array,longtable}
\usepackage{hyperref}
\hypersetup{colorlinks, linkcolor=blue!50!black, urlcolor=blue!50!black}

\lstdefinestyle{gsmd}{
  language=[LaTeX]TeX, basicstyle=\ttfamily\footnotesize, columns=fullflexible,
  keepspaces, commentstyle=\color{gray}\itshape,
  texcsstyle=*\color{blue!60!black},
  moretexcs={abstract,representable,subtype,aggregation,grouping,point,
    attribute,multiattribute,component,generalization,specialization,link,
    note,subdomain,grid,setup,inputgsmdiagram,gsmsetup},
}
\lstset{style=gsmd}
\newtcblisting{example}{enhanced, listing side text, colback=white,
  colframe=black!25, boxrule=0.5pt, arc=1pt, lefthand width=0.5\linewidth,
  listing options={style=gsmd}, halign lower=center, valign lower=center,
  before skip=6pt, after skip=8pt}
\newcommand\cs[1]{\texttt{\textbackslash#1}}
\newcommand\meta[1]{\textlangle\textit{#1}\textrangle}
\newcommand\opt[1]{\texttt{[}#1\texttt{]}}

\title{\textbf{gsmdiagram} --- GSM diagrams for \LaTeX\\[2pt]
  \large Generic Semantic Model, version 3.0.2}
\author{Leunami52}
\date{2026-11-07}

\begin{document}
\maketitle
\begin{center}\small Distributed under the \LaTeX\ Project Public License, version 1.3c or later.\end{center}
\tableofcontents

% =====================================================================
\section{Introduction}

\texttt{gsmdiagram} draws diagrams of the \emph{Generic Semantic Model}
in the usual style (triangles for abstract types, rounded boxes for
representable ones, circles for subtypes). It is
the sister package of \texttt{erdiagram} and works exactly the same way:
a diagram is a list of short, declarative commands inside a
\texttt{gsmdiagram} environment, and the commands exist only inside that
environment.

\paragraph{Installation.} Copy \texttt{gsmdiagram.sty} next to your
document and load it with \verb|\usepackage{gsmdiagram}|. Only TikZ is
needed. The package may be loaded together with \texttt{erdiagram}.
Public names: the environment \texttt{gsmdiagram}, \cs{inputgsmdiagram}
and \cs{gsmsetup}.

\begin{example}
\begin{gsmdiagram}
  \abstract{person}(0,0){Person}
  \representable{pname}(-1.2,2){PName}
  \attribute{person}{pname}{has name}
  \abstract{company}(5,0){Company}
  \attribute{person}{company}{works for}
  \subtype{tourist}(0,-2.4){Tourist}
  \specialization{tourist}{person}
\end{gsmdiagram}
\end{example}

Node commands have the form
\cs{type}\opt{\meta{options}}\verb|{|\meta{id}\verb|}(|\meta{x,y}\verb|){|\meta{label}\verb|}|;
edge commands have the form
\cs{kind}\opt{\meta{options}}\verb|{|\meta{from}\verb|}{|\meta{to}\verb|}{|\meta{label}\verb|}|,
where the label is optional. Options the package does not know are
passed to TikZ.

% =====================================================================
\section{New in version 3.0}

Diagrams from earlier versions keep compiling, and look the same unless they use \texttt{scale}:

\begin{itemize}
\item \texttt{scale=0.5} scales the finished diagram, with its nodes, text and lines. Before, it went
  to TikZ, which moved the positions but not the nodes, so they overlapped.
\item \texttt{sans serif} sets all text of a diagram in the document's sans-serif font.
\end{itemize}

% =====================================================================
\section{The environment and external files}

\verb|\begin{gsmdiagram}[|\meta{options}\verb|] ... \end{gsmdiagram}| accepts the same
options as \texttt{erdiagram}: \texttt{grid} (labelled helper grid),
\texttt{scale} (it scales the whole diagram), \texttt{width}, \texttt{height}, \texttt{max width}, \texttt{max height},
\texttt{checks=false}, any style key of \cs{gsmsetup}, and any
\texttt{tikzpicture} option. Inside, \cs{setup}\verb|{...}| changes style
keys locally and plain TikZ commands may be mixed in.

\verb|\inputgsmdiagram[|\meta{options}\verb|]{|\meta{file}\verb|}| typesets a diagram stored
in a file that contains only the \emph{body} of a \texttt{gsmdiagram}
environment (recommended extension: \texttt{.gsmd}). This is the format
written by the desktop editor. Paths are relative to the main document;
with \texttt{subfiles} use \verb|\inputgsmdiagram{\subfix{file.gsmd}}|.
For files the editor should read back: one command per line, absolute
positions \verb|(x,y)|, options as \texttt{key=value}.

% =====================================================================
\section{Types}

\begin{center}
\verb|\abstract[...]{id}(x,y){Label}|\qquad
\verb|\representable[...]{id}(x,y){Label}|\qquad
\verb|\subtype[...]{id}(x,y){Label}|
\end{center}

Abstract types are triangles, representable types rounded boxes,
subtypes circles. Option \texttt{color=}\meta{colour} sets the
background (a plain colour name is used as a light tint, \texttt{blue}
$\to$ \texttt{blue!20}); TikZ options such as \texttt{minimum width} or
positioning keys work as well. Use \verb|\\| for line breaks.

\begin{example}
\begin{gsmdiagram}
  \abstract{a}(0,0){Person}
  \abstract[color=blue]{b}(3.2,0){Vehicle}
  \representable{c}(0,-1.6){Name}
  \representable[color=yellow]{d}(3.2,-1.6)
    {Zip code}
  \subtype{e}(0,-3.4){Job}
  \subtype[color=green]{f}(3.2,-3.4)
    {Traveling\\salesman}
\end{gsmdiagram}
\end{example}

% =====================================================================
\section{Type constructors: aggregation and grouping}

\begin{center}
\verb|\aggregation[...]{id}(x,y){Label}|\qquad
\verb|\grouping[...]{id}(x,y){Label}|
\end{center}

An aggregation (crossed circle) defines a class composed of other
classes, a grouping (crossed circle with a bar) a set of objects of one
type. The label is printed in red next to the symbol. With the default
\texttt{label pos=auto} it is put where there is the most room
(preferring the right-hand side); \texttt{label pos=north},
\texttt{north east}, \dots\ fix the position. The components are
connected with \cs{component} (dashed arrows).

\begin{example}
\begin{gsmdiagram}
  \grouping{persons}(0,1.2){Persons}
  \aggregation{person}(0,0){Person}
  \component{persons}{person}
  \representable{dob}(2.6,-1){Date of birth}
  \aggregation{name}(-2,-1.2){Name}
  \component{person}{name}
  \component{person}{dob}
  \representable{first}(-2.9,-2.6){First}
  \representable{last}(-1.1,-2.6){Last}
  \component{name}{first}\component{name}{last}
\end{gsmdiagram}
\end{example}

% =====================================================================
\section{Edges}

\begin{longtable}{@{}>{\ttfamily}p{5cm}p{10.8cm}@{}}
\toprule
\normalfont\textbf{Command} & \textbf{Meaning}\\\midrule
\textbackslash attribute & single-valued attribute: arrow\\
\textbackslash multiattribute & multi-valued attribute: double arrow head\\
\textbackslash component & element of a constructed type (from an
  aggregation or grouping): dashed arrow\\
\textbackslash generalization & disjoint IS-A (generalisation): thick
  arrow from the subtype to the supertype\\
\textbackslash specialization & overlapping IS-A (specialisation):
  double-line arrow from the subtype to the supertype\\
\textbackslash link & plain line without an arrow\\
\bottomrule
\end{longtable}

\begin{example}
\begin{gsmdiagram}
  \abstract{p}(0,0){Person}
  \representable{n}(4,1.2){PName}
  \representable{l}(4,-0.6){Language}
  \attribute{p}{n}{has name}
  \multiattribute{p}{l}{speaks}
  \subtype{e}(-1.6,-2.4){Employee}
  \subtype{s}(1.6,-2.4){Student}
  \specialization{e}{p}\specialization{s}{p}
  \abstract{v}(0,-4.6){Vehicle}
  \subtype{c}(-1.6,-6.6){Car}
  \subtype{b}(1.6,-6.6){Boat}
  \generalization{c}{v}\generalization{b}{v}
\end{gsmdiagram}
\end{example}

\subsection{Edge options}

\begin{longtable}{@{}>{\ttfamily}p{4cm}p{11.8cm}@{}}
\toprule
\normalfont\textbf{Option} & \textbf{Meaning}\\\midrule
label side=\meta{side} & \texttt{auto} (default: above flat edges, right
  of steep ones), \texttt{above}, \texttt{below}, \texttt{left},
  \texttt{right}\\
label pos=\meta{fraction} & position along the edge (default 0.5)\\
label segment=\meta{n} & with way points: the segment that carries the
  label (default: the middle one)\\
sloped / horizontal & labels are rotated along their edge by default;
  \texttt{horizontal} prints this one label horizontally
  (\texttt{label sloped=false} does this for a whole diagram)\\
label color=\meta{colour} & colour of the label, e.g.\ for conditions\\
label style=\{\meta{TikZ}\} & any further options for the label\\
bend=\meta{angle} & curved edge; positive bends to the left\\
via=\{(x,y) (x,y) \dots\} & route through way points\\
ortho=-| \textrm{or} |- & right-angled edge\\
\meta{TikZ option} & e.g.\ \texttt{bend left=20}, \texttt{out=90, in=0},
  \texttt{looseness=2}\\
\bottomrule
\end{longtable}

Labels follow their edge by default; \texttt{horizontal} switches a
single label, \texttt{label sloped=false} a whole diagram:

\begin{example}
\begin{gsmdiagram}
  \abstract{p}(0,0){Person}
  \representable{a}(3.5,1.6){Address}
  \representable{b}(3.5,-1.6){Birthday}
  \attribute{p}{a}{lives at}
  \attribute[horizontal]{p}{b}{born on}
\end{gsmdiagram}
\end{example}

\begin{example}
\begin{gsmdiagram}[label sloped=false]
  \abstract{p}(0,0){Person}
  \representable{a}(3.5,1.6){Address}
  \representable{b}(3.5,-1.6){Birthday}
  \attribute{p}{a}{lives at}
  \attribute{p}{b}{born on}
\end{gsmdiagram}
\end{example}

Edges may start or end at a particular anchor of a node
(\verb|person.east|, \verb|person.60|), and \cs{point}\verb|{id}(x,y)|
creates an invisible junction, e.g.\ to merge two edges:

\begin{example}
\begin{gsmdiagram}
  \abstract{p}(0,0){Person}
  \abstract{e}(5,0){Enterprise}
  \attribute{p}{e}{works for}
  \attribute[bend=-25, label side=above]{e}{p}
    {works for$^{-1}$}
  \point{j}(2.5,-1.6)
  \link[ortho=|-]{p.south}{j}
  \link[ortho=|-]{e.south}{j}
  \representable{n}(2.5,-3){Number}
  \attribute[label side=right, horizontal]{j}{n}
    {year(s) of employment}
\end{gsmdiagram}
\end{example}

\begin{example}
\begin{gsmdiagram}
  \abstract{p}(0,0){Person}
  \representable{d}(3.4,0){Date of birth}
  \attribute{p}{d}
  \representable{a}(3.4,-2){Age}
  \attribute[label side=right, horizontal,
    label color=red!70!black,
    label style={font=\scriptsize}]{d}{a}
    {[year(today) $-$ year(d)]}
  \subtype{t}(0,-3){Teenager}
  \specialization[label side=right,
    label color=red!70!black]{t}{p}{[age $<$ 20]}
\end{gsmdiagram}
\end{example}

% =====================================================================
\section{Notes, subdomains and the grid}

These work exactly as in \texttt{erdiagram}:
\verb|\note[|\meta{options}\verb|](x,y){text}| (red free text; \texttt{color=} changes
the colour), \verb|\subdomain[color=..., dotted, fill=...]{Label}{(x,y) (x,y) ...}| (a
frame behind the diagram, two points give a rectangle, more points a
polygon, the label hangs from the first point) and
\verb|\grid(x1,y1)(x2,y2)|.

\begin{example}
\begin{gsmdiagram}
  \subdomain[color=red!60, dotted]{People}
    {(-1.8,1) (1.8,-2.3)}
  \abstract{p}(0,0){Person}
  \subtype{t}(0,-1.5){Tourist}
  \specialization{t}{p}
  \note[color=blue!70!black](1.1,-1.5){[derived]}
  \note[anchor=west](2.1,-0.6)
    {Persons travel\\at least once a year}
\end{gsmdiagram}
\end{example}

% =====================================================================
\section{Rule checks}\label{sec:rules}

Violations of the GSM structure rules are reported as warnings
(\texttt{Package gsmdiagram Warning: GSM rule: \dots}); switch them off
with \verb|\gsmsetup{rule warnings=false}| or the environment option
\texttt{checks=false}.

\begin{itemize}
\item an IS-A edge (\cs{generalization}, \cs{specialization}) that does
  not start at a subtype, or that ends at something other than a subtype
  or an abstract type; an IS-A edge from a node to itself;
\item a component edge that does not start at an aggregation or grouping;
\item a grouping without exactly one component;
\item an aggregation with fewer than two components;
\item a subtype without any IS-A edge to a supertype.
\end{itemize}

The desktop editor checks all of this and more, such as cycles of IS-A edges and the common
root of a subtype's IS-A paths (see its README).

% =====================================================================
\section{Style parameters}

Set globally with \cs{gsmsetup}, per diagram as environment options or
inside a diagram with \cs{setup}.

\begin{longtable}{@{}>{\ttfamily}p{5.6cm}>{\ttfamily}p{3.6cm}p{6.4cm}@{}}
\toprule
\normalfont\textbf{Key} & \normalfont\textbf{Default} & \textbf{Meaning}\\\midrule
\endhead
abstract fill / draw / text & gray!15 / black / black & triangles\\
abstract line width / font & 0.6pt / \textbackslash bfseries & \\
abstract minimum width / height & 2.6cm / 0.6cm & \\
abstract apex angle & 90 & \\
repr fill / draw / text & gray!15 / black / black & rounded boxes\\
repr line width / font & 0.6pt / \textbackslash bfseries & \\
repr minimum width / height & 2.4cm / 0.9cm & \\
repr corner radius & 6pt & \\
sub fill / draw / text & gray!15 / black / black & circles\\
sub line width / font & 0.6pt / \textbackslash bfseries & \\
sub minimum size & 1.4cm & \\
aggr fill / draw & white / black & crossed circles\\
aggr line width / size & 0.6pt / 0.8cm & \\
aggr label color / font & red!70!black / \textbackslash small & \\
aggr label pos & auto & \\
edge color / line width & black / 0.5pt & \\
arrow length / width & 6pt / 4pt & arrow heads\\
generalization line width & 2pt & thick IS-A arrows\\
specialization double distance & 1.5pt & double IS-A arrows\\
label font / color & \textbackslash footnotesize \textbackslash itshape / black & edge labels\\
label pos / side / sloped & 0.5 / auto / true & \\
note color / font & red!70!black / \textbackslash small & \\
subdomain \dots & & as in \texttt{erdiagram}\\
rule warnings & true & \\
sans serif & false & all text of the diagram in the document's sans-serif font\\
grid color / grid label color & gray!25 / gray!55 & \\
\bottomrule
\end{longtable}

{\raggedright TikZ styles for deeper changes: \texttt{gsm abstract}, \texttt{gsm repr},
\texttt{gsm subtype}, \texttt{gsm constructor},
\texttt{gsm constructor label}, \texttt{gsm label}, \texttt{gsm edge},
\texttt{gsm single}, \texttt{gsm multi}, \texttt{gsm component},
\texttt{gsm generalization}, \texttt{gsm specialization},
\texttt{gsm note}.\par}

% =====================================================================
\section{A complete example: a music festival}

The sample diagram of ERDTeX (\emph{File $\to$ Open Example}): abstract and
representable types, subtypes with generalizations and a specialization, an
aggregation and a grouping, single- and multi-valued attributes, a derived one,
and two attribute lines in opposite directions.

\noindent\inputgsmdiagram[width=\linewidth]{festival.gsmd}

\lstinputlisting[basicstyle=\ttfamily\scriptsize, frame=single,
  rulecolor=\color{black!25}]{festival.gsmd}

% =====================================================================
\section{Known limitations}

\begin{enumerate}
\item \textbf{Edge labels do not avoid other elements.} Labels are
  placed beside their edge, but collisions with other edges or nodes
  are not detected; adjust \texttt{label side}, \texttt{label pos} or
  the layout.
\begin{example}
\begin{gsmdiagram}
  \abstract{a}(0,0){A}
  \representable{b}(4,0){B}
  \representable{c}(4,-1){C}
  \attribute{a}{b}{a long label}
  \attribute{a}{c}{another long label}
\end{gsmdiagram}
\end{example}
\item \textbf{Arrow heads on double lines.} With strongly curved
  specialisation edges the double line may look uneven near the head.
\item \textbf{Cycles} of IS-A edges and the ``common root'' rule are not
  checked here (only an edge from a node to itself is reported); the desktop editor does check them.
\item \textbf{Way points} in \texttt{via} must be plain coordinates or
  node anchors.
\end{enumerate}

\end{document}
